En esta etapa, (\texttt{app/services/prediccion/features.py}), se adapta el historial de demanda a un formato que contenga las variables necesarias que los modelos usan como entrada (\textit{features}):

\begin{itemize}
    \item Componentes de calendario
    \item Rezagos de la demanda
    \item Estadísticas móviles
    \item Estadísticas históricas por producto
\end{itemize}

La implementación de esta etapa empieza con un merge incial, en este merge se trae las categorías abc y variabilidad a la tabla. Aunque no se usan como variables de entrada (\texttt{feature\_cols}), se incluyen porque posteriormente se necesitaran para mostrarlas en el dashboard y tabla de predicciones:

\Needspace{9\baselineskip}
\captionof{listing}{Unión de las categorías ABC y de variabilidad a la tabla mensual.}\label{lst:features_merge}
\begin{minted}{python}
def construir_features(df_mensual, producto_comportamiento):
    # ...
    df_mensual = df_mensual.merge(
        producto_comportamiento[["producto_id", "categoria_abc", "variabilidad"]],
        on="producto_id", how="left"
    )

\end{minted}

Posteriormente, se extraen tanto el año como el mes de cada ejemplo, esto adaptado del repositorio \textit{Demand Forecasting \& Inventory Optimization for Consumer Products} \cite{ejemplo1_demandforecast}, al que en este trabajo se le referira como \textit{Ejemplo 1 (SKU Demand Forecasting)}. Con la diferencia de que en Thary, estos datos se incluyeron con el fin de en la implementación que se tomó como guia, el mes "month", queda como un feature final para el modelo, mientras que en Thary, este dato posteriormente sea tranformado en \texttt{mes\_sin} y \texttt{mes\_cos}, proceso que se explicara más adelante:

\Needspace{4\baselineskip}
\captionof{listing}{Extracción del año y del mes de cada registro.}\label{lst:features_anio_mes}
\begin{minted}{python}
    df_mensual["año"] = df_mensual["fecha"].dt.year
    df_mensual["mes"] = df_mensual["fecha"].dt.month
\end{minted}

Un rezago (lag), se usa en Machine Learning y hace referencia al valor anterior a una variable, en este caso, representa el valor de la demanda del mes anterior, el cual es usado como dato de entrada para predecir el mes actual. Los lags junto con las medias y desviaciones móviles se encargan de capturar el comportamiento reciente de la demanda de cada producto. En Thary, para cada producto se tiene:

\begin{itemize}
    \item lag\_1: demanda del mes anterior
    \item lag\_2: demanda de hace dos meses
    \item lag\_3: demanda de hace tres meses
\end{itemize}

El modelo los usa como features porque la demanda pasada ayuda a poder determinar cuanto se va a vender el proximo mes. Por eso es que un producto necesita mínimo 3 meses de historial antes de los primeros 3 meses del dataset para poder calcular lag1, lag2 y lag3. El problema era que el sistema solo conoce los datos que vienen dentro del mismo archivo CSV que el usuario sube, por lo que no tiene forma de saber que paso antes del primer mes que aparece en el archivo.

Si un dataset empieza en enero de 2021, el lag1 de enero necesitaría la demanda de diciembre de 2020, dato que no esta en el historial dado por el usuario, solo seria a partir de abril 2021 que el sistema contaría con los 3 lags comparados correctamente. Es debido a que no hay información suficiente para calcular los lags correctamente, que el sistema descarta los primeros meses a la hora de entrenar el modelo mediante (\texttt{dropna()}. Como esos primeros 3 meses se descartan,  el valor de mes de esos meses nunca aparece como algo que el modelo tuvo que predecir durante el entrenamiento, por lo que estaría el riesgo de que, si más adelante se le pide al modelo predecir justo ese mismo mes del año siguiente (por ejemplo, enero 2022), el modelo no tenga ninguna referencia previa de que número de mes es ese, y termine adivinando mal, generando predicciones muy bajas o poco confiables para ese mes en particular

Para evitar este problema, se decidió que durante esta etapa de features, se representaría el mes utilizando seno y coseno en vez de usar directamente el número del mes (1-12). Esto permite que el modelo interprete el mes como un punto continuo cercano a los meses vecinos que si fueron entrenados, por ejemplo, diciembre y enero quedan numericamente uno cerca del otro, lo cual tiene mucho más sentido que hacer que el sistema los vea como el número 12 y 1, valores mucho más lejanos de lo que deberían ser para dos meses vecinos. 

Una vez que el modelo entienda la cercania real de los meses, es capaz de dejar de interpretarlos como categorías independientes, y puede generar una predicción razonable de un mes que nunca halla visto antes (enero, por ejemlo), en base la información de un mes cercano (diciembre).

A la hora de armar la función para construir los features se tomaron en cuenta diferentes implementaciones previas de distintos estudios con diferentes propuestas para un sistema de predicción de demanda, se compararon y se seleccionaron aquellas que se apataban mejor a las necesidades del sistema y le aportaban valos al mismo, junto con la implementación de nuevas ideas y mejoras que se fueron propuestas e implemantadas en Thary. Una de estas propuestas nuevas fue la codificación ciclica del mes, que como ya se explicó antes, consiste en representar los meses mediante dos valores calculados con seno y coseno, como se puede ver a continuación:

\Needspace{5\baselineskip}
\captionof{listing}{Codificación cíclica del mes mediante seno y coseno.}\label{lst:features_ciclico}
\begin{minted}{python}
    # ...
    df_mensual["mes_sin"] = np.sin(2 * np.pi * df_mensual["mes"] / 12)
    df_mensual["mes_cos"] = np.cos(2 * np.pi * df_mensual["mes"] / 12)
\end{minted}

Por otra parte, la metodologia de implementar el uso de lags, también se tomó del \textit{Ejemplo 1 (SKU Demand Forecasting)} \cite{ejemplo1_demandforecast} y se adaptó a las necesidades del sistema Thary, agregando la codificación cíclica de los meses y la inclusión de features adicionales como el inventario y el precio.

\Needspace{8\baselineskip}
\captionof{listing}{Cálculo de los rezagos \texttt{lag\_1}, \texttt{lag\_2}, \texttt{lag\_3} y \texttt{lag\_6} por producto.}\label{lst:features_lags}
\begin{minted}{python}
    # ... 
    # (Ejemplo 1 - SKU Demand Forecasting)
    df_mensual["lag_1"] = df_mensual.groupby("producto_id")["demanda"].shift(1)
    df_mensual["lag_2"] = df_mensual.groupby("producto_id")["demanda"].shift(2)
    df_mensual["lag_3"] = df_mensual.groupby("producto_id")["demanda"].shift(3)
    df_mensual["lag_6"] = df_mensual.groupby("producto_id")["demanda"].shift(6)
\end{minted}

Otra estrategia que se tomó del \textit{Ejemplo 1 (SKU Demand Forecasting)} \cite{ejemplo1_demandforecast}, fue el uso de las estadísticas móviles. Estas aportan al sistema ya que consisten en un resumen de varios meses juntos. Mientras que los rezagos solo representan cuanto se vendió exactamente en un mes, las estadísticas móviles son un resumen de varios meses juntos.

\begin{itemize}
\item \texttt{media\_movil\_3} y \texttt{media\_movil\_6}: Estas variables representan cual es el nivel reciente de la demanda de los productos calculando el promedio de los últimos 3 y 6 meses. Esto ayuda que el modelo pueda distingir un producto que vende de manera estable de uno irregular, ya que de ser ventas regulares, la media móvil de 3 meses será muy similar a la de 6 meses, mientras que si es inestable, ambas medias serán muy diferentes.
\item \texttt{desviacion\_movil\_3}: Representa que tan errática fue la demanda reciente, es decir, de los últimos 3 meses. Los productos que cuentan con demanda estable tendrán una desviación baja mientras que los productos con ventas inestable tendrán una desviación alta.
\end{itemize}

En general, el modelo puede ver el comportamiento general reciente de los productos más allá de solo los lags.

En cuanto a la implementación, se incluyó en la función de (\texttt{construir\_features}), en el que cada línea calcula por producto (\texttt{groupby("producto\_id")}) una estadística sobre una ventana de meses anteriores (\texttt{.shift(1)}) para no incluir el mes actual, evitando fuga de datos.


\Needspace{10\baselineskip}
\captionof{listing}{Cálculo de las estadísticas móviles \texttt{media\_movil\_3}, \texttt{media\_movil\_6} y \texttt{desviacion\_movil\_3}.}\label{lst:features_moviles}
\begin{minted}{python}
    # ...
    # (Ejemplo 1 - SKU Demand Forecasting)

    df_mensual["media_movil_3"] = df_mensual.groupby("producto_id")["demanda"].transform(
        lambda x: x.shift(1).rolling(3).mean()
    )
    df_mensual["media_movil_6"] = df_mensual.groupby("producto_id")["demanda"].transform(
        lambda x: x.shift(1).rolling(6).mean()
    )
    df_mensual["desviacion_movil_3"] = df_mensual.groupby("producto_id")["demanda"].transform(
        lambda x: x.shift(1).rolling(3).std()
    )
\end{minted}

Por otra parte, para la construcción de la etapa de Features también se tomó en cuenta el repositorio en Kaggle \textit{Forecasting Inventory Demand Using XGBoost} \cite{ejemplo3_ksmooi_inventorydemand} o \textit{Ejemplo 3 (XGBoost)} \cite{ejemplo3_ksmooi_inventorydemand}. De este repositorio, se tomó el uso de el promedio del producto, pero esta vez, calculados en base a todo el historial completo y no solo en base a los últimos 3 o 6 meses. 

Como extensión propuesta por el proyecto, junto al promedio también se incluyó la desviación del producto, esto para poder seguir una misma estructura en todo el sistema, ya que si se tiene la (\texttt{media\_movil\_3}) y la (\texttt{desviacion\_movil\_3}), es decir, el nivel y la variabilidad de los últimos 3 meses, es conveniente tener la media y desviación de todo el historial del producto.

Estos valores ya fueron calculados en la etapa de segmentación, y simplemente son agregados en los features para que el modelo también los use como variables de entrada.

\Needspace{10\baselineskip}
\captionof{listing}{Incorporación de la media y la desviación históricas de cada producto.}\label{lst:features_historicas}
\begin{minted}{python}
    # ...
    # (Forecasting Inventory Demand Using XGBoost)
    stats_cols = ["demanda_total", "media", "desviacion"]
    df_mensual = df_mensual.merge(
        producto_comportamiento[["producto_id"] + stats_cols].rename(
            columns={"media": "media_producto", "desviacion": "desviacion_producto"}
        ),
        on="producto_id", how="left"
    )
\end{minted}

Después, se calculca cuanto stock habia el mes anterior. La condición (\texttt{if "inventario" in df\_mensual.columns:}) determina si el csv que fue subido por el usuario contiene la columna "inventario", esta es una columna opcional que simplemente representa cuantas unidades de ese producto (stock disponible) habia fisicamente en el inventario en ese mes, Thary funciona sin ella, pero si el usuario la incluye, el modelo la usara como una variable de entrada más para poder brindarle un análisis mucho más detallado. 

(\texttt{(inventario\_lag\_1)}) es el inventario del mes anterior, con (\texttt{.groupby("producto\_id")}) se asegura que este se calcule por separado para cada producto, esto con la intención de evitar que el inventario de un producto se vaya a mezclar con otro, esto usando (\texttt{(shift(1))}) para poder obtener el inventario del mes anterior. Si no se hace esto, el modelo podría aprender que el inventario de un producto es igual al del mes actual, lo cual no es correcto y generaría predicciones erroneas. 

\Needspace{6\baselineskip}
\captionof{listing}{Cálculo del inventario del mes anterior (\texttt{inventario\_lag\_1}).}\label{lst:features_inventario_lag}
\begin{minted}{python}
    # ...
    # (Forecasting Inventory Demand Using XGBoost))
    if "inventario" in df_mensual.columns:
        df_mensual["inventario_lag_1"] = df_mensual.groupby("producto_id")["inventario"].shift(1)
\end{minted}

En esta parte, se implementó otra idea del \textit{Ejemplo 1 (SKU Demand Forecasting)} \cite{ejemplo1_demandforecast} para poder descartar las filas sin historial suficiente para tener los lags. Para esto, se define una lista con nueve nombres de las columnas obligatorias para poder entrenar, con (\texttt{(dropna(subset=feature\_history))}) se revisan las 4 columnas y elimina cualquier fila que tenga un espacio vacío, posteriormente el resultado se guarda en (\texttt{(df\_model)}), que es la tabla final que se usara para entrenar el modelo.

\Needspace{10\baselineskip}
\captionof{listing}{Descarte de las filas sin historial suficiente para calcular los rezagos.}\label{lst:features_dropna}
\begin{minted}{python}
    # ...
    # (Inventory Demand Forecast (XGBoost Reg))
    feature_history = ["lag_1", "lag_2", "lag_3", "media_movil_3"]
    df_model = df_mensual.dropna(subset=feature_history).copy()

    if len(df_model) == 0:
        raise ValueError(
            "Ningún producto tiene al menos 3 meses seguidos de historial, que es lo mínimo "
            "para calcular tendencias. Sube un archivo con más meses continuos de datos."
        )
\end{minted}

Otro de los conceptos que fueron propuestos en Thary es el tratamiento que se decidió darle a la variable con los precios de venta de los productos. Aunque el CSV pueda incluir esa variable, el modelo solo usa su valor actual (el precio de ese mes), sin calcular cómo cambió respecto al mes anterior. Esto debido a que por naturaleza, el precio de un producto puede influir en la demanda de un producto, ya que cuando algo baja de precio, la gente tiende a comprar más, mientras que cuando este sube, las personas compran menos. El problema es que si esta regla se repite muy seguido, se corre el riesgo de que el modelo aprenda que el cambio de precio predice muy bien la demanda, entonces, aunque el precio es una variable que si se tiene en cuenta, se excluye por completo su variabilidad de un mes a otro.  

\Needspace{10\baselineskip}
\captionof{listing}{Lista de variables de entrada (\texttt{feature\_cols}) y variables opcionales de precio e inventario.}\label{lst:features_lista}
\begin{minted}{python}
    # ...
    feature_cols = [
        "año", "mes_sin", "mes_cos",
        "lag_1", "lag_2", "lag_3", "lag_6",
        "media_movil_3", "media_movil_6", "desviacion_movil_3",
        "media_producto", "desviacion_producto"
    ]
    if "precio" in df_model.columns:
        feature_cols.append("precio")
    if "inventario" in df_model.columns:
        feature_cols.append("inventario")
        if "inventario_lag_1" in df_model.columns:
            feature_cols.append("inventario_lag_1")
\end{minted}

Finalmente, como extensión propia de Thary, se hace un relleno final. Para esto, se recorre cada una de las columnas que van a usar los modelos, si hay alguna fila vavia en la tabla, se reemplazan los vacios de la columna con 0 mediante (\texttt{df\_model[col] = df\_model[col].fillna(0)}). Esto es importante porque si el modelo recibe un producto que solo aparece a partir de cierto mes, por ejemplo, abril 2021 en un dataset que empieza en enero de ese mismo año, el mes de abril pudo pasar el filtro de dropna() porque ya tiene (\texttt{(lag\_1, lag\_2, lag\_3 y media\_movil\_3 completos)}, pero en este caso, el lag\_6 necesitaría la demanda de octubre de 2020, mes que no existe en el archivo porque el producto todavua no tiene 6 meses de historial.

La fila pasaría el filtro porque \texttt{lag\_6} no está en \texttt{feature\_history}, pero inevitablemente terminará en \texttt{feature\_cols}, y ahí es donde se ejecuta este bloque final, si \texttt{lag\_6} es vacío en esa fila se rellena con 0.

\Needspace{6\baselineskip}
\captionof{listing}{Relleno final de los valores vacíos de las variables de entrada con 0.}\label{lst:features_relleno}
\begin{minted}{python}
    # ...
    for col in feature_cols:
        if df_model[col].isna().any():
            df_model[col] = df_model[col].fillna(0)
\end{minted}
